%------------------------------------------------------------------------------%
\section{Exercícios}

\begin{exercises*}

%------------------------------------------------------------------------------%
\begin{exercise}
  Nos Teoremas e proposições a seguir, destaque a hipótese e a tese.
  \begin{alphalist}
    \item Se $a,b$ e $c$ são lados de um triângulo e $\hat{A}$ o ângulo
      oposto ao lado $a$, então $a^2 = b^2 + c^2 - 2bc \cos \hat{A}$
    \item Toda função diferenciável em $x=a$ é contínua em $x=a$. (Ou de
      modo equivalente: Se $f$ é uma função diferenciável então $f$ é uma
      função contínua em $x=a$.)
    \item Se $T$ é um triângulo equilátero então $Q$ é um triângulo
      isósceles (Ou de modo equivalente : Todo triângulo equilátero é
      isósceles).
    \item Se $x=1$, então $x^2=1$
  \end{alphalist}

  \begin{answer}
    \begin{alphalist}
      \item Hipótese: $a$, $b$ e $c$ são lados de um triângulo e $\hat{A}$
            é o ângulo oposto ao lado $a$ \\
            Tese: $a^2 = b^2 + c^2 - 2bc \cos \hat{A}$
      \item Hipótese: $f$ é uma função diferenciável em $x=a$ \\
            Tese: $f$ é uma função contínua em $x=a$
      \item Hipótese: $T$ é um triângulo equilátero \\
            Tese: $T$ é um triângulo isósceles
      \item Hipótese: $x=1$ \\
            Tese: $x^2=1$
    \end{alphalist}
  \end{answer}
\end{exercise}

%------------------------------------------------------------------------------%
\begin{exercise}
  Para as sentenças a seguir, dê a recíproca e a contrapositiva das
  proposições que seguem:
  \begin{alphalist}
    \item Se $k$ é um número par primo então $k=2$.
    \item Toda função diferenciável em $x=a$ é contínua em $x=a$
    \item Se $T$ é um triângulo equilátero então $T$ é um triângulo
      isósceles (Ou de modo equivalente : Todo triângulo equilátero é
      isósceles).
    \item Se $x=1$, então $x^2=1$
  \end{alphalist}

  \begin{answer}
    \begin{alphalist}
      \item Recíproca: Se $k=2$, então $k$ é um número par primo \\
        Contrapositiva: Se $k \neq 2$, então $k$ não é um número par primo
      \item Recíproca: Se $f$ é contínua em $x=a$,
        então $f$ é diferenciável em $x=a$ \\
        Contrapositiva: Se $f$ não é contínua em $x=a$,
        então $f$ não é diferenciável em $x=a$
      \item Recíproca: Se $T$ é isósceles, então $T$ é equilátero \\
        Contrapositiva: Se $T$ não é isósceles, então $T$ não é equilátero
      \item Recíproca: Se $x^2=1$, então $x=1$ \\
        Contrapositiva: Se $x^2 \neq 1$, então $x \neq 1$
    \end{alphalist}
  \end{answer}
\end{exercise}

%------------------------------------------------------------------------------%
\begin{exercise}
  A recíproca da proposição apresentada na questão anterior, item $e$ é
  verdadeira? Justifique.

  \begin{answer}
    Não. A recíproca afirma que se $x^2=1$ então $x=1$.
    Contudo, se $x = -1$, temos que $x^2 = 1$.
    Assim,  $x=\pm 1 \neq 1$, portanto a recíproca é falsa
  \end{answer}
\end{exercise}


%------------------------------------------------------------------------------%
\begin{exercise}
  Demonstre, de duas maneiras diferentes, as seguintes proposições. Em
  cada item, você deve escolher dois tipos de demonstração para fazer
  dentre: demonstração direta, redução ao absurdo e usando a
  contrapositiva.
  \begin{alphalist}
    \item A soma de dois números pares é sempre um número par.
    \item A soma de dois números ímpares é sempre um número par.
    \item O produto de dois números ímpares é sempre um número ímpar.
    \item Se $x^2$ é ímpar, então $x$ é ímpar.
    \item Se $x+y$ é divisível por $2$ então $x-y$ é divisível por $2$.
  \end{alphalist}

  \begin{answer}
    \begin{alphalist}
        \item
            Demonstração Direta:
            Sejam $n = 2k$ e $m = 2j$.
            A soma é $n + m = 2k + 2j = 2(k + j)$, que é par \\
            Redução ao Absurdo: Suponha que $2k + 2j = 2p + 1$.
            Então $2(k + j - p) = 1$, o que é uma contradição (par igual a ímpar)

        \item
            Demonstração Direta:
            Sejam $n = 2k + 1$ e $m = 2j + 1$.
            A soma é $(2k + 1) + (2j + 1) = 2(k + j + 1)$, que é par \\
            Redução ao Absurdo:
            Suponha que $2k+1 + 2j+1 = 2p + 1$.
            Então $2(k + j - p +1) = 1$, o que é uma contradição (par igual a ímpar)

        \item
            Demonstração Direta:
            Sejam $n = 2k + 1$ e $m = 2j + 1$.
            O produto é $4kj + 2k + 2j + 1 = 2(2kj + k + j) + 1$, que é ímpar \\
            Redução ao Absurdo:
            Suponha que o produto seja par ($2p$).
            Então $2(2kj+k+j) + 1 = 2p \Rightarrow 1 = 2(p - 2kj - k - j)$,
            o que é uma contradição (par igual a ímpar)

        \item
            Usando a Contrapositiva:
            Se $x=2k$, então $x^2 = 4k^2 = 2(2k^2)$, que é par.
            Provada a contrapositiva, a original é verdadeira \\
            Redução ao Absurdo:
            Suponha que $x$ é par.
            Então $x^2 = (2k)^2 = 2(2k^2)$ que é par,
            o que contradiz a hipótese de $x^2$ ser ímpar

        \item
            Demonstração Direta:
            Se $x + y = 2k$, então $x - y = (x + y) - 2y = 2k - 2y = 2(k - y)$,
            que é divisível por 2 \\
            Redução ao Absurdo:
            Suponha que $x - y$ seja ímpar ($2j+1$).
            Somando com $x+y=2k$, teríamos $2x = 2(k+j)+1$,
            o que é uma contradição (par igual a ímpar)
    \end{alphalist}
  \end{answer}
\end{exercise}


%------------------------------------------------------------------------------%
\begin{exercise}
  Demonstre, usando redução ao absurdo, que $\sqrt{3}$ é um número
  irracional.

  \begin{answer}
    Redução ao absurdo: Suponha $\sqrt{3} = \frac{p}{q}$ irredutível.
    Então $3 = \frac{p^2}{q^2} \Rightarrow p^2 = 3q^2$.
    Logo, $p = 3k$ \\
    Substituindo: $(3k)^2 = 3q^2 \Rightarrow 3k^2 = q^2$.
    Logo, $q$ também é múltiplo de 3, contradizendo a irredutibilidade
  \end{answer}
\end{exercise}


%------------------------------------------------------------------------------%
\begin{exercise}
  Todas as proposições abaixo são falsas. Dê um contra-exemplo que
  invalide cada uma das afirmações.
  \begin{alphalist}
    \item Em todo triângulo, de lados $a,b$ e $c$, e com $a$ sendo o maior
      lado, vale que $a^2=b^2+c^2$.
    \item Se $a$ e $b$ são números não-negativos então $\sqrt{a+b}=
      \sqrt{a}+ \sqrt{b}$.
    \item Se $a$ é um número real qualquer, então $a^2>a$.
    \item Todo número natural ímpar é primo.
    \item Todo triângulo é retângulo.
  \end{alphalist}

  \begin{answer}
    \begin{alphalist}
      \item Contraexemplo: Triângulo de lados $a=5, b=4, c=2$.
            Temos $5^2 \neq 4^2 + 2^2$
      \item Contraexemplo: $a=9, b=16$.
            Temos $\sqrt{9+16}=5$, mas $\sqrt{9}+\sqrt{16}=7$
      \item Contraexemplo: $a=0.5$.
            Temos $0.5^2 = 0.25$, e $0.25 < 0.5$
      \item Contraexemplo: $n=9$ é ímpar,
            mas não é primo (é divisível por 3)
      \item Contraexemplo: Triângulo equilátero possui todos os
            ângulos iguais a $60^\circ$
    \end{alphalist}
  \end{answer}
\end{exercise}

\end{exercises*}

%------------------------------------------------------------------------------%
